\documentclass[12pt,a4paper]{article}

% ---------- پکیج‌ها (همه قبل از xepersian) ----------
\usepackage{fontspec}
\usepackage{amsmath, amssymb, amstext}
\usepackage{graphicx}
\usepackage{geometry}
\usepackage{xcolor}
\usepackage{fancyhdr}

% ---------- تنظیم حاشیه ----------
\geometry{margin=1in}

% ---------- تنظیم fancyhdr ----------
\pagestyle{fancy}
\setlength{\headheight}{15pt} % حل خطای headheight

% هدرها
\lhead{کوییز درس مکانیک‌تحلیلی}
\chead{دانشکده فیزیک دانشگاه صنعتی خواجه نصیرالدین طوسی}
\rhead{\today}

% ---------- بسته فارسی و فونت ----------
\usepackage{xepersian}
\settextfont{B Nazanin}

\begin{document}
	
	% جدول اطلاعات فردی
	\begin{tabular}{|p{7cm}|p{7cm}|}
		\hline
		\textbf{نام و نام خانوادگی:} & \textbf{شماره دانشجویی:} \\
		\hline
	\end{tabular}
	
	% بخش سوالات
	\section*{سوالات:}
	\begin{enumerate}
		\item 
		رابطه بین مکان و زمان این متحرک به صورت زیر است:
		\[
		t = a x^3 - b x + c
		\]
		که در آن \(a\)، \(b\) و \(c\) مقادیر ثابتی هستند.
		\begin{itemize}
			\item الف) سرعت این متحرک را بدست آورید.
			\item ب) شتاب این متحرک را بدست آورید.
			\item ج) دیمانسیون هر یک از ثوابت \(a\)، \(b\) و \(c\) را بدست آورده و با دیمانسیون سرعت و شتاب مقایسه کنید.
		\end{itemize}
		
		\vspace{0.8cm} % فاصله بین سوال 1 و 2
		
		% سوال ۲
		\item 
		یک دستگاه مختصات سه‌بعدی \((x, y, z)\) را در نظر بگیرید.
		\begin{itemize}
			\item الف) ماتریس دوران حول محور \(z\) به زاویه‌ی \(\theta\) را بنویسید و نشان دهید که این ماتریس یک ماتریس متعامد است.
			\item ب) ماتریس کلی دوران حول محور \(y\) به زاویه‌ی \(\phi\) را بنویسید.
			\item ج) ترکیب دو دوران متوالی را در نظر بگیرید: ابتدا دوران به اندازه‌ی \(\theta\) حول محور \(z\)، سپس دوران به اندازه‌ی \(\phi\) حول محور \(y\). ماتریس کلی این ترکیب را به‌دست آورید (به صورت حاصل‌ضرب دو ماتریس دوران).
			\item د) نشان دهید که دترمینان ماتریس دوران کلی برابر با \(+1\) است.
			\item هـ) بررسی کنید که آیا ترتیب دوران‌ها (اول حول \(z\) سپس حول \(y\)) با حالت معکوس آن یکی است یا نه.
		\end{itemize}
	\end{enumerate}
	
	\vspace{10pt}
	\textbf{موفق باشید.}
	
\end{document}
